\chapter{Sourcing Tools and AI-Assisted Search}
\label{ch:sourcing-tools}

Finding the right person and their contact details is a mechanical bottleneck in every outreach method this
guide covers. This chapter is a directory of the specific tools cited across the source material, plus two
AI-driven search workflows that automate parts of the job hunt itself.

\section{Email and Contact-Finding Tools}

\begin{longtable}{p{0.2\textwidth} p{0.76\textwidth}}
\toprule
\textbf{Tool} & \textbf{What it does, and caveats} \\
\midrule
\endhead
RocketReach & Input a LinkedIn profile URL to fetch a professional work email. Reportedly does not consume
credits when no contact information is found. \\
ContactOut & Search by name, job title, and company; can surface work emails, personal emails, and phone
numbers. \\
Skrapp.io & Verifies email addresses by checking validity and server status before you send. \\
Apollo.io & Browser extension — open a LinkedIn profile and click ``access email address'' to retrieve business
emails and company phone numbers where available. Roughly a 50/50 success rate per profile, per the source. Also
useful for exploring a company's professional email addresses and inferring its address format
(\texttt{firstname.lastname@company.com}, etc.) — but a guessed address is not a verified one. \\
Signal Hire & Browser extension for personal contact info when a professional email isn't available via other
tools. Limited to roughly 10 credits/month on the cited plan. \\
Hunter.io & Enter a company name to explore associated professional email addresses and infer the organization's
common address format. Also has a Chrome extension for finding a name and email directly from a company
website. \\
Clearbit (Connect) & Another email-discovery option, used similarly to Hunter. \\
Prospeo & Browser extension, used alongside Apollo for finding a recruiter's professional email from their
LinkedIn profile. \\
\bottomrule
\end{longtable}

\begin{caution}
Before emailing anyone found this way: confirm the address is actually associated with the right person and
company; prefer contact details clearly intended for business communication; never send mass unsolicited
messages or repeatedly contact someone who has asked not to be contacted; and don't use personal contact details
obtained through inappropriate means. The point of these tools is to make \emph{relevant} outreach possible, not
to bypass someone's reasonable privacy expectations.
\end{caution}

\section{Scraping Techniques for Building Larger Lead Lists}

Two manual scraping workflows cited, both built from free or low-cost tools chained together:

\begin{description}
  \item[YouTube description scraping] Use a Google advanced-search operator —
  \texttt{site:youtube.com "[keyword]" "[location]" "gmail.com"} — to surface channels that list a contact
  email in their description. Install an Email Extractor Chrome extension to automatically detect and collect
  valid addresses from the resulting pages, then export the list to a spreadsheet.
  \item[LinkedIn/Twitter scraping via Google search] Use a query-generation site (cited: \emph{Recruiting}) to
  build a structured search string from a social network, job title, and location. Refine it in
  \textbf{WordCounter.net} by adding a target email provider in quotes (\texttt{"gmail.com"}, \texttt{"yahoo.com"},
  etc.) to narrow results to a specific domain. Run the search on Google, copy the full results page, and paste
  it into \textbf{Email Extractor Lite} to strip everything down to a clean list of addresses. The source notes
  roughly 30--40 minutes of work per batch for a meaningfully sized list, and suggests varying job titles,
  locations, and email providers (AOL, Outlook, Hotmail) to widen coverage.
\end{description}

\begin{caution}
These techniques surface publicly listed contact information at scale — the same etiquette rules from the
contact-finding tools above apply even more strongly here: relevance first, respect opt-outs, and don't treat a
large list as license to spam it indiscriminately.
\end{caution}

\section{Qualifying Leads and Running Infrastructure at Scale}

Two tools cited for the operational side of larger campaigns (full workflow in
Chapter~\ref{ch:case-study-emails}):

\begin{description}
  \item[Meta Ads Library] Search by keyword to find brands currently running paid social ads — active ad spend
  is a plausible (not certain) signal that a company also has budget for other marketing services.
  \item[Semrush] Check a prospect's apparent business age, search visibility, and branded search traffic to
  separate plausibly-established businesses from very new or low-visibility ones before spending outreach effort
  on them.
  \item[Instantly] A cold-email platform for managing contact lists, scheduling multi-step sequences, tracking
  opens, and coordinating multiple sending mailboxes across domains.
\end{description}

\section{AI Agents for Automated Job Search}

Two distinct workflows use an AI assistant with browser/connector access to search job boards directly, rather
than you searching manually.

\subsection{Claude + Naukri: The ``Under 50 Applicants'' Strategy}

The core idea: most roles on a crowded job portal are saturated within the first wave of applicants, so
targeting \emph{recently posted, lightly applied-to} listings meaningfully improves your odds of being seen.
Setup: install the Claude Chrome extension, keep your LinkedIn profile open as a reference, and enable the
extension's permissions in Claude's settings.

\begin{template}
\textbf{Basic prompt}\\[4pt]
Use my LinkedIn profile as a reference. Analyse my skills, experience, and certifications. Then go to Naukri
portal and find job openings for [YOUR TARGET ROLE] in [YOUR TARGET CITY OR INDIA]. Make sure the roles match my
LinkedIn skill set. Give me the direct Naukri links so I can apply for each one.
\end{template}

\begin{template}
\textbf{Advanced prompt (with the under-50-applicants filter)}\\[4pt]
Use my LinkedIn profile as a reference. Analyse my complete skill set, experience, and certifications. Then go
to Naukri portal and find job openings for [Product Manager] in [INDIA] that: were posted in the last 24 hours
only, and have fewer than 50 applicants at the time of search. Make sure each listing genuinely matches my
LinkedIn profile keywords and skills. Return the direct Naukri application links for each matched role.
\end{template}

\subsection{Claude + Job-Board Connectors (Indeed and Others)}

A second workflow uses Claude's free-tier connectors (an Indeed connector, specifically) to search and filter
listings by tech stack, salary, and location, including a comparative view to help decide between matches. The
same approach extends to freelance and contract roles, and to other connectors (ZipRecruiter, Workable, Dice) for
broader coverage.

\begin{template}
\textbf{Remote job prompt}\\[4pt]
Find high-quality full-time remote Frontend Engineer jobs on Indeed.\\[2pt]
Focus on: React.js, Next.js, TypeScript, remote-first companies, companies hiring internationally.\\[2pt]
Target country: [Country Name]. Industry: [Fintech / SaaS / E-commerce / AI / etc.]\\[2pt]
Salary range — minimum: [\$30k+ / AED 15k+ / 10 Lakh etc.]\\[2pt]
Prioritize: recently posted jobs, less than 100 applicants, product-based companies, good engineering culture.\\[2pt]
Exclude: onsite roles, hybrid roles, staffing agencies.
\end{template}

\begin{template}
\textbf{Freelancing prompt}\\[4pt]
Find freelance, contract, or client-based frontend development opportunities on Indeed.\\[2pt]
Focus on: React.js, Next.js, TypeScript, SaaS dashboards, web apps, landing pages.\\[2pt]
Target country: [Country Name]. Industry: [Fintech / SaaS / E-commerce / AI / etc.]\\[2pt]
Budget/contract preference: high-paying clients only, long-term contracts preferred.\\[2pt]
Prioritize: startup founders hiring directly, remote freelance opportunities, fast-moving startups.\\[2pt]
Exclude: low-budget gigs, cheap outsourcing agencies, non-remote opportunities.
\end{template}

\begin{keytakeaway}
Both prompt templates are reusable: adapt the role, stack, country, industry, and salary line to your situation,
and treat the ``exclude'' block as seriously as the ``focus'' block — it's what keeps the AI from burying good
matches under a pile of irrelevant ones.
\end{keytakeaway}

\section{Choosing a Tool: Quick Reference}

\begin{longtable}{p{0.34\textwidth} p{0.6\textwidth}}
\toprule
\textbf{You need to...} & \textbf{Reach for...} \\
\midrule
\endhead
Find a specific person's email from their LinkedIn & Apollo.io, Prospeo, RocketReach, Hunter.io \\
Verify an email before sending & Skrapp.io \\
Find a personal email when the professional one is missing & Signal Hire \\
Build a large lead list from public descriptions/profiles & Google search operators + Email Extractor
(Chrome) \\
Check whether a company has marketing budget/activity & Meta Ads Library, Semrush \\
Run a multi-step outreach sequence at scale & Instantly \\
Search a job board against your own profile automatically & Claude + Naukri or an Indeed/job-board connector \\
\bottomrule
\end{longtable}

